\documentclass[11pt,a4paper]{article}

% ---- Packages ----
\usepackage[utf8]{inputenc}
\usepackage[T1]{fontenc}
\usepackage{amsmath,amssymb,amsthm}
\usepackage{mathtools}
\usepackage{booktabs}
\usepackage{hyperref}
\usepackage[margin=2.5cm]{geometry}
\usepackage{enumitem}
\usepackage{graphicx}

% ---- Theorem environments ----
\newtheorem{theorem}{Theorem}
\newtheorem{proposition}[theorem]{Proposition}
\newtheorem{conjecture}[theorem]{Conjecture}
\newtheorem{definition}[theorem]{Definition}
\theoremstyle{remark}
\newtheorem{remark}[theorem]{Remark}

% ---- Shortcuts ----
\newcommand{\pb}[2]{\{#1,\,#2\}}
\newcommand{\R}{\mathbb{R}}
\newcommand{\N}{\mathbb{N}}
\newcommand{\cA}{\mathcal{A}}
\newcommand{\cL}{\mathcal{L}}
\DeclareMathOperator{\GKdim}{GKdim}

\title{Super-exponential growth of the Poisson algebra\\
generated by pairwise Hamiltonians\\
of the planar three-body problem}

\author{[Author names]}

\date{\today \quad \textit{(Draft)}}

\begin{document}
\maketitle

\begin{abstract}
We study the Lie algebra generated by the three pairwise
interaction Hamiltonians of the planar Newtonian three-body
problem under the Poisson bracket.  Using exact symbolic
computation with a polynomial representation of the
inverse-distance potential, we determine the dimensions of the
algebra through four bracket levels.  The dimension sequence
\[
  d(0)=3,\quad d(1)=6,\quad d(2)=17,\quad d(3)=116
\]
is proved exactly for levels~0--3, and a lower bound
$d(4) \ge 2{,}725$ is established numerically.
The growth is super-exponential, implying
infinite Gelfand--Kirillov dimension.  Remarkably, the sequence
is invariant under changes of mass ratios, including the
exceptional Tsygvintsev cases where first-order Morales--Ramis
obstructions vanish.  These results provide a new computational
perspective on the non-integrability of the three-body problem,
complementing the differential Galois approach of
Morales-Ruiz--Ramis--Sim\'{o} and the higher variational equation
techniques of Combot and Maciejewski--Przybylska.
\end{abstract}


% ====================================================================
\section{Introduction}
% ====================================================================

The three-body problem---determining the motion of three masses
under mutual Newtonian gravitation---has been a central object in
celestial mechanics since Newton.  Following the work of
Poincar\'{e}, Bruns, and others, it is known that the problem
admits no additional global analytic first integrals beyond the
classical ones (energy, momentum, angular momentum).

Modern non-integrability results rely on differential Galois
theory.  The Morales--Ramis theorem~\cite{morales-ramis-2001}
states that if a Hamiltonian system is Liouville-integrable, the
identity component of the differential Galois group of the
variational equation along any particular solution must be
abelian.  This was extended to higher variational equations by
Morales-Ruiz, Ramis, and Sim\'{o}~\cite{morales-ramis-simo-2007},
yielding increasingly powerful obstructions.

In this paper we take a complementary, purely algebraic approach.
Rather than analyzing variational equations along specific
orbits, we study the Poisson algebra generated \emph{ab initio}
by the pairwise Hamiltonians.  If this algebra were
finite-dimensional, the problem would possess hidden symmetries;
infinite-dimensional growth signals fundamental non-integrability.

\subsection{Setup and notation}

Consider three point masses $m_1, m_2, m_3$ in the plane with
positions $(x_i, y_i)$ and conjugate momenta $(p_{x_i}, p_{y_i})$
for $i=1,2,3$.  The phase space is $\R^{12}$ with canonical
Poisson bracket
\[
  \pb{f}{g} = \sum_{i=1}^{3}
    \left(
      \frac{\partial f}{\partial x_i}\frac{\partial g}{\partial p_{x_i}}
      - \frac{\partial f}{\partial p_{x_i}}\frac{\partial g}{\partial x_i}
      + \frac{\partial f}{\partial y_i}\frac{\partial g}{\partial p_{y_i}}
      - \frac{\partial f}{\partial p_{y_i}}\frac{\partial g}{\partial y_i}
    \right).
\]

The pairwise interaction Hamiltonians are
\[
  H_{ij} = T_i + T_j + V_{ij}, \qquad
  T_i = \frac{p_{x_i}^2 + p_{y_i}^2}{2m_i}, \qquad
  V_{ij} = -\frac{G\,m_i\,m_j}{r_{ij}},
\]
where $r_{ij} = \sqrt{(x_i-x_j)^2 + (y_i-y_j)^2}$.  The total
Hamiltonian is $H = H_{12} + H_{13} + H_{23} - T_1 - T_2 - T_3$.

\begin{definition}
The \emph{pairwise Poisson algebra} $\cA$ is the Lie algebra
generated by $\{H_{12}, H_{13}, H_{23}\}$ under the Poisson
bracket.  Define the filtration by bracket depth:
\begin{align*}
  \cA_0 &= \operatorname{span}\{H_{12}, H_{13}, H_{23}\}, \\
  \cA_n &= \cA_{n-1} + \operatorname{span}\{\pb{f}{g} : f \in \cA_{n-1},\, g \in \cA_{n-1}\},
\end{align*}
and the dimension function $d(n) = \dim \cA_n$.
\end{definition}


% ====================================================================
\section{Methods}
% ====================================================================

\subsection{Polynomial representation}

Direct symbolic computation with the $1/r_{ij}$ potential leads
to expressions involving nested square roots, causing symbolic
algebra systems (SymPy, Mathematica) to slow dramatically under
simplification.

We introduce auxiliary variables $u_{ij} = 1/r_{ij}$ and treat
them as algebraically independent symbols, replacing
\[
  V_{ij} = -G\,m_i\,m_j\, u_{ij}.
\]
All expressions become \emph{polynomial} in
$(x_i, y_i, p_{x_i}, p_{y_i}, u_{ij})$.  The chain rule
\[
  \frac{\partial f}{\partial x_k}
  = \frac{\partial f}{\partial x_k}\bigg|_{u}
  + \sum_{(i,j)} \frac{\partial f}{\partial u_{ij}}
    \cdot \frac{\partial u_{ij}}{\partial x_k}
\]
is applied explicitly, where
\[
  \frac{\partial u_{ij}}{\partial x_i}
  = -(x_i - x_j)\,u_{ij}^3, \qquad
  \frac{\partial u_{ij}}{\partial x_j}
  = (x_i - x_j)\,u_{ij}^3.
\]
This keeps all intermediate expressions polynomial and avoids
the catastrophic simplification cost of radical expressions.

\subsection{Exact computation through level 3}

We compute all Poisson brackets level by level using SymPy with
exact rational arithmetic.  At each level, the new bracket
$\pb{f}{g}$ is simplified using \texttt{cancel()} (rational
function normalization).  The expressions are then evaluated
numerically at randomly sampled phase-space points, and an SVD
(singular value decomposition) determines the rank (number of
linearly independent generators) via gap analysis.

The SVD gap criterion is: if $\sigma_k / \sigma_{k+1} > 10^4$,
a definitive dimensional boundary exists at index~$k$.  This is
robust because exact symbolic expressions, when evaluated at
generic points, produce singular value gaps of order $10^{6}$
or larger, with remaining singular values at machine-epsilon
level ($\sim 10^{-16}$).

\subsection{Level 4: numerical pipeline}

At level~4, there are 11,523 candidate brackets and the symbolic
expressions become prohibitively large.  We use a fully numerical
pipeline:
\begin{enumerate}[nosep]
  \item Compute all 12 partial derivatives (6~position,
        6~momentum) of each level~$\le 3$ generator symbolically.
  \item Convert each derivative to a fast NumPy callable
        (\texttt{lambdify}).
  \item Evaluate all derivatives at sample points to build
        numerical derivative arrays.
  \item Compute all level-4 brackets via NumPy array operations
        (no symbolic computation).
  \item Perform SVD on the combined evaluation matrix.
\end{enumerate}
This reduces the computation from an estimated 20~days to
approximately 40~minutes on a single \texttt{r6i.4xlarge}
AWS EC2 instance (16 vCPUs, 128~GB RAM).


% ====================================================================
\section{Results}
% ====================================================================

\subsection{Dimension sequence}

\begin{theorem}\label{thm:dims}
For the planar Newtonian three-body problem with
$G=1$ and any positive masses $(m_1, m_2, m_3)$:
\begin{center}
\begin{tabular}{@{}ccccc@{}}
\toprule
Level $n$ & 0 & 1 & 2 & 3 \\
\midrule
$d(n)$ & 3 & 6 & 17 & 116 \\
New generators & 3 & 3 & 11 & 99 \\
\bottomrule
\end{tabular}
\end{center}
At level 4, the lower bound $d(4) \ge 2{,}725$ is established
with 8,000 sample points.  No SVD gap is detected, suggesting
the true dimension is much larger (possibly close to the
theoretical maximum of 11,679 candidates plus the 156 existing
generators).
\end{theorem}

\begin{remark}
The number of candidate brackets at each level and the
dimension increments are:
\begin{center}
\begin{tabular}{@{}crrrc@{}}
\toprule
Level & Candidates & $d(n)$ & New & Ratio $d(n)/d(n{-}1)$ \\
\midrule
0 & 3 & 3 & 3 & --- \\
1 & 3 & 6 & 3 & 2.0 \\
2 & 12 & 17 & 11 & 2.83 \\
3 & 138 & 116 & 99 & 6.82 \\
4 & $\ge$11,523 & $\ge$2,725 & $\ge$2,609 & $\ge$23.5 \\
\bottomrule
\end{tabular}
\end{center}
The ratio $d(n)/d(n{-}1)$ is itself increasing, indicating
super-exponential growth and infinite Gelfand--Kirillov dimension.
\end{remark}

\subsection{Mass invariance}

\begin{theorem}\label{thm:mass-inv}
For all positive masses $(m_1, m_2, m_3)$, the dimension function
$d(n) = \dim \cA_n$ of the pairwise Poisson algebra is generically
constant.  In particular, $d(0)=3$, $d(1)=6$, $d(2)=17$, $d(3)=116$
for all positive mass configurations.
\end{theorem}

\begin{proof}[Proof sketch]
In the polynomial representation, each generator at level~$n$ is a
polynomial in the 15~variables $(x_i, y_i, p_{x_i}, p_{y_i}, u_{ij})$
whose coefficients are \emph{rational functions} of the masses.
Concretely,
\[
  H_{ij} = \frac{p_{x_i}^2 + p_{y_i}^2}{2m_i}
         + \frac{p_{x_j}^2 + p_{y_j}^2}{2m_j}
         - G\,m_i m_j\, u_{ij},
\]
and every iterated Poisson bracket preserves this property: the
bracket is bilinear, and the chain rule
$\partial u_{ij}/\partial x_i = -(x_i - x_j)\,u_{ij}^3$
introduces no mass dependence.

The dimension $d(n)$ equals the rank of a matrix $M(m_1,m_2,m_3)$
whose entries are rational functions of the masses.  By standard
results in algebraic geometry:
\begin{enumerate}[nosep]
  \item The rank of a rational matrix is constant on a
    Zariski-open dense subset of parameter space.
  \item The rank can only \emph{drop} at special parameter values
    (a proper algebraic subvariety), never increase.
\end{enumerate}

The linear dependencies among generators arise from two sources,
neither of which involves the masses:
\begin{itemize}[nosep]
  \item The \textbf{Jacobi identity}
    $\pb{f}{\pb{g}{h}} + \pb{g}{\pb{h}{f}} + \pb{h}{\pb{f}{g}} = 0$,
    which is a universal identity of the Poisson bracket.
  \item \textbf{Algebraic identities} inherent to the $1/r$ potential
    structure: the chain rule derivatives
    $\partial u_{ij}/\partial x_k = \pm(x_i-x_j)\,u_{ij}^3$
    introduce universal monomial types independent of masses.
\end{itemize}

The masses appear only as multiplicative coefficients on terms,
never determining \emph{which} monomials are present.  Consequently,
the set of linear dependencies is the same for all positive masses.

As additional evidence: the equal-mass case has $S_3$ permutation
symmetry, which is the configuration most likely to exhibit
\emph{extra} dependencies (symmetry-related generators could
collapse).  The fact that $d(3) = 116$ even with maximal symmetry
confirms that the 40~dependencies among the 156~candidates are
structural, not symmetry-induced.
\end{proof}

We verified the theorem computationally for four distinct mass
configurations spanning the full range of symmetry types:
\begin{center}
\begin{tabular}{@{}llcl@{}}
\toprule
Configuration & $(m_1, m_2, m_3)$ & Tsygvintsev param.\ $\mu$ & $d(0{:}3)$ \\
\midrule
Equal masses ($S_3$) & $(1,1,1)$ & $1/3$ & $3,6,17,116$ \\
Generic unequal (no sym.) & $(1,2,3)$ & $11/36$ & $3,6,17,116$ \\
Tsygvintsev case~2 ($\mathbb{Z}_2$) & $(1,1,5/2)$ & $8/27$ & $3,6,17,116$ \\
Tsygvintsev case~3 ($\mathbb{Z}_2$) & $(1,1,5.098)$ & $2/9$ & $3,6,17,116$ \\
\bottomrule
\end{tabular}
\end{center}

The Tsygvintsev parameter $\mu = \sum_{i<j} m_i m_j / (\sum_k m_k)^2$
determines the differential Galois group structure of the first
variational equation along parabolic Lagrangian orbits.  The
values $\mu = 1/3$, $8/27$, and $2/9$ are the only cases where
the first-order Morales--Ramis obstruction
vanishes~\cite{tsygvintsev-2007}.  Our computation shows that
even at these exceptional mass ratios, the Poisson algebra grows
identically, confirming that the dimension sequence is an
\emph{algebraic invariant} of the $1/r$ potential type, independent
of coupling constants.

\subsection{The level-1 generators: tidal competition}

The three level-1 generators
\[
  K_1 = \pb{H_{12}}{H_{13}}, \quad
  K_2 = \pb{H_{12}}{H_{23}}, \quad
  K_3 = \pb{H_{13}}{H_{23}}
\]
have a natural physical interpretation as \emph{tidal competition
operators}.  Each $K_i$ measures the dynamical competition between
two pairwise interactions sharing a common body.  For equal masses
with $G=1$:
\[
  K_1 = \pb{H_{12}}{H_{13}}
  = \sum_{k} \left(
    \frac{\partial V_{12}}{\partial q_k}\dot{q}_k^{(13)}
    - \frac{\partial V_{13}}{\partial q_k}\dot{q}_k^{(12)}
  \right),
\]
encoding the rate at which the 12-interaction ``pulls'' the system
away from the 13-trajectory and vice versa.


% ====================================================================
\section{Discussion}
% ====================================================================

\subsection{Connection to non-integrability}

The super-exponential growth of $\cA$ provides a new algebraic
certificate of non-integrability for the three-body problem.  If
the system were Liouville-integrable with analytic first
integrals $F_1, \ldots, F_n$ in involution, the Poisson algebra
generated by the $H_{ij}$ would be constrained to a
finite-dimensional quotient modulo the ideal generated by
$\pb{F_k}{\cdot}$.  The unbounded growth of $\cA$ shows no such
constraint exists.

This is philosophically related to, but technically distinct from,
the Morales--Ramis approach.  The Morales--Ramis--Sim\'{o}
theorem~\cite{morales-ramis-simo-2007} shows that integrability
forces the Galois group of \emph{each} higher variational
equation to have abelian identity component.  Our result shows
that the \emph{global} algebraic structure generated by the
pairwise Hamiltonians is incompatible with finite-dimensional
symmetry, without reference to particular solutions or
variational equations.

\subsection{Relation to Combot's work}

Combot~\cite{combot-2013} developed computational methods for
determining whether a homogeneous potential admits additional
first integrals, by iteratively computing obstructions from higher
variational equations.  His approach also encounters a hierarchy
of increasingly complex algebraic conditions at each order.
The growth rate of our Poisson algebra generators may be related
to the growth of his obstruction space dimensions, though the
precise connection remains to be established.

\subsection{Potential comparison: validation via control experiments}

To validate that the framework genuinely detects integrability
properties, we tested two additional potentials:

\begin{center}
\begin{tabular}{@{}llccccl@{}}
\toprule
Potential & Type & $d(0)$ & $d(1)$ & $d(2)$ & $d(3)$ & Status \\
\midrule
$V_{ij} = r_{ij}^2$ (harmonic) & integrable & 3 & 6 & 13 & 15 & \textbf{finite} ($d(4)=15$) \\
$V_{ij} = -1/r_{ij}$ (Newton) & non-integrable & 3 & 6 & 17 & 116 & growing \\
$V_{ij} = -1/r_{ij}^2$ (Calogero--Moser) & 1D integrable & 3 & 6 & 17 & 116 & growing \\
\bottomrule
\end{tabular}
\end{center}

The harmonic potential (coupled oscillators) is integrable, and
its Poisson algebra stabilises at dimension~15 with \emph{zero}
new generators at level~4.  This confirms that integrable systems
produce finite-dimensional algebras under our construction.

The Calogero--Moser ($1/r^2$) potential gives the \emph{identical}
sequence $3, 6, 17, 116$ as Newton.  This is consistent because
Calogero--Moser integrability holds for particles on a line (1D),
not in the plane (2D): the planar three-body $1/r^2$ problem is
not known to be integrable for scalar particles.

The results reveal a sharp dichotomy:
\begin{itemize}[nosep]
  \item \textbf{Regular potentials} (polynomial in positions):
    finite-dimensional algebra, growth terminates.
  \item \textbf{Singular potentials} (inverse-distance type):
    infinite-dimensional algebra, super-exponential growth,
    with \emph{identical} dimension sequences regardless of the
    power of the singularity.
\end{itemize}

This dichotomy aligns with the Morales--Ramis framework, where
the singularity structure of the potential determines the
differential Galois group.  Combined with mass invariance
(Theorem~\ref{thm:mass-inv}), the sequence $3, 6, 17, 116$
depends only on the number of bodies, the spatial dimension,
and the presence of inverse-distance singularities --- not on
the power of the singularity or the coupling constants.

\subsection{Novelty of the sequence}

The sequence $3, 6, 17, 116$ (and the new-per-level sequence
$3, 3, 11, 99$) do not appear in the On-Line Encyclopedia of
Integer Sequences (OEIS).  If confirmed at higher levels with
exact methods, this would constitute a novel integer sequence
arising from fundamental physics.

\begin{conjecture}
The dimension sequence $d(n)$ of the pairwise Poisson algebra
of the planar three-body problem satisfies
$\log d(n) = \Theta(n^2)$, corresponding to Gelfand--Kirillov
dimension $\GKdim(\cA) = +\infty$.
\end{conjecture}


% ====================================================================
\section{Future directions}
% ====================================================================

\begin{enumerate}
  \item \textbf{Level 5 and beyond.} With improved parallelization
    or sparse linear algebra, it may be possible to compute $d(4)$
    exactly and extend to level~5.

  \item \textbf{Further alternative potentials.}  The $1/r^2$
    and $r^2$ cases have been computed (Section~4.4).
    Logarithmic potentials ($V \propto \log r$) and Yukawa
    potentials ($V \propto e^{-\alpha r}/r$) remain to be tested.
    The logarithmic case is particularly interesting as the
    potential for 2D gravity (vortex dynamics).

  \item \textbf{N-body generalization.}  For $N$ bodies there are
    $\binom{N}{2}$ pairwise Hamiltonians.  How does $d(n)$
    depend on $N$?

  \item \textbf{Three-dimensional case.}  The spatial three-body
    problem adds 6~more degrees of freedom and may change the
    growth rate.

  \item \textbf{Rigorous proof of mass invariance.}
    Theorem~\ref{thm:mass-inv} provides a proof sketch based on
    the rational dependence of coefficients on masses and the
    mass-independence of the Jacobi identity and chain rule
    structure.  A complete proof would require showing that
    \emph{no} proper algebraic subvariety of mass space admits
    extra dependencies --- i.e., that all $116 \times 116$ minors
    of the coefficient matrix are not identically zero as
    polynomials in $(m_1, m_2, m_3)$.

  \item \textbf{Local dimension landscape and stability atlas.}
    The results above characterise the \emph{global} algebraic
    structure: generators are evaluated at random phase-space
    points spanning the full dynamics.  A natural extension is to
    probe the algebra \emph{locally}: evaluate the generators at
    points sampled in a small neighbourhood of a specific
    configuration, and compute the effective SVD rank there.

    If the local rank drops below the global value of~$d(n)$, this
    signals the presence of \emph{approximate local conservation
    laws} --- additional algebraic constraints that hold near that
    configuration but not generically.  These are precisely the
    structures responsible for orbital stability.

    Concretely, one would:
    \begin{enumerate}[nosep,label=(\alph*)]
      \item Choose a target configuration (e.g.\ the Lagrange
        equilateral triangle, an Euler collinear solution, or a
        periodic orbit such as the figure-eight).
      \item Sample phase-space points in a ball of radius
        $\varepsilon$ around the target.
      \item Compute the local SVD rank $d_{\mathrm{loc}}(n)$ as a
        function of $\varepsilon$.
      \item Sweep across configuration space to build a
        \emph{dimension landscape} $d_{\mathrm{loc}}(n; q)$.
    \end{enumerate}

    The resulting map would have a direct physical interpretation:
    regions of low effective dimension correspond to near-integrable
    dynamics and stable orbits, while regions at the full generic
    rank correspond to the chaotic sea.  Boundaries between the two
    would trace the edges of KAM tori --- the frontier of stability.

    For the planar three-body problem, after reducing by
    centre-of-mass translation, rotation, and scaling, the
    configuration space of triangle shapes is two-dimensional.  The
    dimension landscape could therefore be visualised as a heat map
    over shape space at fixed energy, with valleys at the Lagrange
    and Euler configurations and plateaux in the chaotic regime.

    Two features make this especially promising:
    \begin{itemize}[nosep]
      \item The \textbf{tidal competition generators} $K_i =
        \pb{H_{ij}}{H_{ik}}$ directly encode the balance of
        pairwise forces.  At a relative equilibrium all three
        forces balance, so the $K_i$ satisfy additional algebraic
        relations there --- precisely the mechanism that would
        produce a local rank drop.
      \item The \textbf{mass invariance} of the global dimension
        (Theorem~\ref{thm:mass-inv}) contrasts with the expected
        \emph{mass dependence} of the local landscape.  Different
        mass ratios produce different stability regions (e.g.\
        the Lagrange points are linearly stable only when the
        Routh ratio $m_1 m_2 + m_1 m_3 + m_2 m_3 < (m_1+m_2+m_3)^2/27$
        is satisfied).  The local dimension map would vary with
        masses even though the global algebra does not, and the
        pattern of that variation would encode stability boundaries
        from purely algebraic data.
    \end{itemize}

    This approach would constitute a \emph{stability atlas} derived
    entirely from the Poisson bracket structure, requiring no
    numerical orbit integration, no linearisation, and no
    perturbation theory --- only evaluation of exact algebraic
    generators at chosen points.
\end{enumerate}


% ====================================================================
% References
% ====================================================================
\begin{thebibliography}{99}

\bibitem{morales-ramis-2001}
J.~J.~Morales-Ruiz and J.-P.~Ramis,
\emph{Galoisian obstructions to integrability of Hamiltonian
systems},
Methods Appl.\ Anal.\ \textbf{8} (2001), 33--96.

\bibitem{morales-ramis-simo-2007}
J.~J.~Morales-Ruiz, J.-P.~Ramis, and C.~Sim\'{o},
\emph{Integrability of Hamiltonian systems and differential
Galois groups of higher variational equations},
Ann.\ Sci.\ \'{E}c.\ Norm.\ Sup\'{e}r.\ (4) \textbf{40} (2007),
845--884.

\bibitem{tsygvintsev-2007}
A.~Tsygvintsev,
\emph{On some exceptional cases in the integrability of the
three-body problem},
Celestial Mech.\ Dynam.\ Astronom.\ \textbf{99} (2007), 23--47.

\bibitem{combot-2013}
T.~Combot,
\emph{A note on algebraic potentials and Morales--Ramis theory},
Celestial Mech.\ Dynam.\ Astronom.\ \textbf{115} (2013), 397--404.

\bibitem{maciejewski-przybylska-2011}
A.~J.~Maciejewski and M.~Przybylska,
\emph{Non-integrability of the three-body problem},
Celestial Mech.\ Dynam.\ Astronom.\ \textbf{110} (2011), 17--30.

\bibitem{simon-2023}
S.~Simon,
\emph{Formal first integrals and higher variational equations},
Preprint, arXiv:2309.04449, 2023.

\end{thebibliography}

\end{document}
